% TOP_PROBLEM: {"id": 8000013, "title": "Perturbation theory for exactly solvable combinatorial problems", "category_id": 16, "category": {"id": 16, "name": "physics", "display_name": "Mathematical Physics", "description": "Problems at the intersection of mathematics and physics.", "slug": "physics", "order_index": 16, "created_at": "2026-07-31T15:26:25.671Z"}, "status": "solved", "statusQualification": "The inherited catalogue status is solved; this historical statement is retained without independently verifying that settlement."}
% FIRST_POSTED: 2026-10-11
% ENTRY_KIND: statement_only
% SUBMISSION: {"kind": "statement_only", "category": "top_problems", "problem_id": "8000013", "proposal_type": "established", "source": "Percy Deift, Some open problems in random matrix theory and the theory of integrable systems (2007), arXiv:0712.0849, Problem 13. https://arxiv.org/abs/0712.0849", "status_evidence": "Inherited source-catalogue status solved preserved. Primary source accessed 11 October 2026; compared against the supplied catalogue excerpt. Historical wording and partial-result remarks were checked for scope; no comprehensive current-literature or proof-verification audit was performed.", "status_checked_on": "2026-10-11", "duplicate_check": "Full 15,773-record source registry checked for identical normalized statements; available site registry checked for identical normalized title/statement and canonical-ID collision. Canonical ID is absent from the site registry. Related canonical records 8100003. No mathematical-equivalence literature audit performed.", "attribution": "Percy Deift. Statement preparation: OpenAI GPT-6 Codex assistant, generation 11 October 2026; source-matching and notation editing by the assistant, no independent human mathematical verification. Inherited rights/provenance note: Public arXiv source TeX; arXiv reuse license NEEDS_REVIEW"}
% SOURCE_STATEMENT: "The asymptotic behavior of a variety of combinatorial problems has now been analyzed in great detail. Here we have in mind Ulam's problem for the length of the longest increasing subsequence, the tiling problem for the Aztec diamond, the hexagon tiling problem, and the last passage percolation problem, amongst many others. In all cases, in an appropriate scaling limit, the statistical fluctuations in the systems at hand are described by random matrix theory. The asymptotic analysis, however, depends in a critical, and rigid way, on the underlying probability measures for the systems. For example, the analysis of the last passage percolation problem requires the waiting time $w_i$ at each site $i$ to be either geometrically or exponentially distributed (Johansson): if the statistics of the waiting times is neither geometric nor exponential, the analysis fails completely. What happens if the geometric distribution, say, is slightly perturbed? The challenge here is to develop an effective perturbation theory for such systems. One expects that the random matrix behavior of the fluctuations should persist. Related work on this problem has been done by Baik-Suidan and Bouchard-Martin (cf. Problem 10)."
% Preparation and final self-audit: CONTRIBUTING.md prompts C and D; source/audit artifacts retained in the local physics_statements_2026_10_11 batch.
% Imported from: physics_statements.tex; UnsolvedMath ID 8000013
% !TeX program = lualatex
\documentclass[11pt]{article}
\usepackage[margin=25mm]{geometry}
\usepackage{fontspec}
\setmainfont{Latin Modern Roman}
\usepackage{amsmath,amssymb,mathtools}
\usepackage{unicode-math}
\setmathfont{Latin Modern Math}
\usepackage{xurl,hyperref}
\hypersetup{colorlinks=true,urlcolor=blue}
\setcounter{secnumdepth}{0}
\setlength{\parindent}{0pt}
\setlength{\parskip}{5pt}
\setlength{\emergencystretch}{3em}
\newcommand{\sref}[2]{\hyperlink{source:#1:#2}{[S#2]}}
\newcommand{\cataloguescope}[1]{\paragraph{Recorded target.}#1}
\newcommand{\R}{\mathbb R}
\newcommand{\C}{\mathbb C}
\newcommand{\N}{\mathbb N}
\newcommand{\Z}{\mathbb Z}
\newcommand{\Q}{\mathbb Q}
\newcommand{\F}{\mathbb F}
\newcommand{\eps}{\varepsilon}
\begin{document}
% UnsolvedMath ID 8000013; source catalogue code AMR-079-0013
\setcounter{section}{8000012}
\section{Perturbation theory for exactly solvable combinatorial problems}\label{8000013}
{\small UnsolvedMath ID 8000013 · Mathematical Physics}
\subsection{Definitions and mathematical statement}

\paragraph{Shared notation.}$\mathbb N=\{1,2,\ldots\}$; the standard number fields have their usual meanings. All additional conventions are stated locally.


\paragraph{Definitions and notation.}Exactly solvable probabilistic models admit algebraic formulae for their distributions. In last-passage percolation, geometric or exponential site-weight laws lead to tractable formulae. A perturbation changes the weight distribution; a universality result would show persistence of a normalized fluctuation limit for a declared class of such changes.

\paragraph{Catalogue wording retained for traceability.}
The asymptotic behavior of a variety of combinatorial problems has now been analyzed in great detail. Here we have in mind Ulam's problem for the length of the longest increasing subsequence, the tiling problem for the Aztec diamond, the hexagon tiling problem, and the last passage percolation problem, amongst many others. In all cases, in an appropriate scaling limit, the statistical fluctuations in the systems at hand are described by random matrix theory. The asymptotic analysis, however, depends in a critical, and rigid way, on the underlying probability measures for the systems. For example, the analysis of the last passage percolation problem requires the waiting time $w_i$ at each site $i$ to be either geometrically or exponentially distributed (Johansson): if the statistics of the waiting times is neither geometric nor exponential, the analysis fails completely. What happens if the geometric distribution, say, is slightly perturbed? The challenge here is to develop an effective perturbation theory for such systems. One expects that the random matrix behavior of the fluctuations should persist. Related work on this problem has been done by Baik-Suidan and Bouchard-Martin (cf. Problem 10).
\paragraph{Scope and source qualification.}The inherited catalogue status is solved; this historical statement is retained without independently verifying that settlement.
\subsection{Short English statement}
Build an effective perturbation theory that preserves random matrix fluctuation behavior beyond exactly solvable probability laws.
\subsection{Sources}
\begin{itemize}
\item[\hypertarget{source:8000013:1}{S1}] ULAM.AI and the UnsolvedMath Contributors, UnsolvedMath, record 8000013 (AMR-079-0013). Catalogue curation under CC BY 4.0; definitions and English exposition added locally.
\url{https://huggingface.co/datasets/ulamai/UnsolvedMath}.
\item[\hypertarget{source:8000013:2}{S2}] Percy Deift, Some open problems in random matrix theory and the theory of integrable systems (2007), arXiv:0712.0849, Problem 13.
\url{https://arxiv.org/abs/0712.0849}.
\end{itemize}
\label{end:8000013}
\end{document}
